%分類モデルごとの各表の説明
%|-テストデータごとの表の説明
% |-表から読み取れること
% |-考察

% \noindent{\red{Todo:左LiDARの6組について、テストデータからStayをなしにした3分類での精度表とマトリクスを示し、精度が上がる/下がるかを考察する。}}\\

% NOTE:
% - Tables are generated into tex/evaluation/confusion-matrices/table.tex

% Confusion matrices (text only; tables are in confusion-matrices-tables.tex)

Tables~\ref{normalnonematrix} and \ref{normalkmeansmatrix} present \deleted{confusion matrices for each subject} \revised{representative confusion matrices for test subjects A to AB}.
%
The rows represent the ground truth labels, and the columns represent the inferred labels.
%
\deleted{Table~\ref{normalnonematrix} and Table~\ref{normalkmeansmatrix} correspond to the results shown in Table~\ref{accresults}(a) and Table~\ref{accresults}(i), respectively.}
%
\begin{revisedblock}
We compare two conditions to examine the effect of clustering.
  %
  The conditions are: normals without clustering and normals with k-means.
\end{revisedblock}
%
\deleted{Both of these used normals as features; however, Table~\ref{accresults}(a) exhibited high accuracy, while Table~\ref{accresults}(i) showed lower accuracy.}
\revised{Both of the used normals as features; however, the normals without clustering tended to show higher accuracy, while the normals with k-means tended to show lower accuracy in Table~\ref{accresults}.}
%
In Table~\ref{normalnonematrix}, the diagonal elements from the top-left to the bottom-right are large, indicating that the inferred labels largely agreed with the ground truth. 
%
However, there were frequent misclassifications of typing as pad operation. 
%
This is likely due to the difficulty in distinguishing whether both hands were touching the pad or the keyboard. 
%
In Table~\ref{normalkmeansmatrix}, in addition to misclassifying typing as pad operation, there were also frequent misclassifications of pad operation as typing.
  
  
% NOTE: subfig's default (a)...(z) overflows for 28 subfloats in a table
% (uses the `subtable` counter). Use numeric labels locally.
{
\renewcommand{\thesubtable}{\arabic{subtable}}
% Auto-generated confusion matrices for subjects A--AB (representative conditions)
% Regenerate: powershell -ExecutionPolicy Bypass -File tex/evaluation/generate_confusion_matrices_A_AB.ps1

\begin{table}[b]
  \centering
  \caption{Matrix of test results with normals and without clustering method. Diagonal toward lower right corner is correct. Vertical columns represent number of labels inferred as mouse operation, pad operation, sitting still, and typing, respectively, from left to right. Horizontal rows represent, from top, number of labels where actual label is mouse operation, pad operation, sitting still, or typing.}
  \label{normalnonematrix}
  \subfloat[Test data A]{
    \label{nnA}
    \scalebox{0.67}{
      \begin{tabular}{|c|c|c|c|}
      \hline
        $273$ & $0$ & $0$ & $0$ \\ \cline{1-4}
        $172$ & $99$ & $2$ & $1$ \\ \cline{1-4}
        $1$ & $0$ & $95$ & $0$ \\ \cline{1-4}
        $17$ & $103$ & $25$ & $131$ \\ \cline{1-4}
      \hline
    \end{tabular}
    }
  }
  \hfill
  \subfloat[Test data B]{
    \label{nnB}
    \scalebox{0.67}{
      \begin{tabular}{|c|c|c|c|}
      \hline
        $272$ & $3$ & $2$ & $0$ \\ \cline{1-4}
        $0$ & $265$ & $3$ & $7$ \\ \cline{1-4}
        $0$ & $1$ & $92$ & $0$ \\ \cline{1-4}
        $0$ & $2$ & $0$ & $274$ \\ \cline{1-4}
      \hline
    \end{tabular}
    }
  }
  \hfill
  \subfloat[Test data C]{
    \label{nnC}
    \scalebox{0.67}{
      \begin{tabular}{|c|c|c|c|}
      \hline
        $279$ & $0$ & $0$ & $0$ \\ \cline{1-4}
        $0$ & $0$ & $0$ & $278$ \\ \cline{1-4}
        $0$ & $0$ & $94$ & $0$ \\ \cline{1-4}
        $0$ & $0$ & $0$ & $276$ \\ \cline{1-4}
      \hline
    \end{tabular}
    }
  }
  \hfill
  \subfloat[Test data D]{
    \label{nnD}
    \scalebox{0.67}{
      \begin{tabular}{|c|c|c|c|}
      \hline
        $267$ & $9$ & $1$ & $0$ \\ \cline{1-4}
        $48$ & $185$ & $43$ & $3$ \\ \cline{1-4}
        $0$ & $6$ & $89$ & $0$ \\ \cline{1-4}
        $0$ & $252$ & $1$ & $22$ \\ \cline{1-4}
      \hline
    \end{tabular}
    }
  }
  \hfill
  \subfloat[Test data E]{
    \label{nnE}
    \scalebox{0.67}{
      \begin{tabular}{|c|c|c|c|}
      \hline
        $275$ & $1$ & $0$ & $2$ \\ \cline{1-4}
        $0$ & $275$ & $2$ & $0$ \\ \cline{1-4}
        $0$ & $63$ & $34$ & $0$ \\ \cline{1-4}
        $0$ & $272$ & $1$ & $4$ \\ \cline{1-4}
      \hline
    \end{tabular}
    }
  }
  \hfill
  \subfloat[Test data F]{
    \label{nnF}
    \scalebox{0.67}{
      \begin{tabular}{|c|c|c|c|}
      \hline
        $276$ & $2$ & $2$ & $0$ \\ \cline{1-4}
        $0$ & $270$ & $5$ & $4$ \\ \cline{1-4}
        $2$ & $4$ & $90$ & $1$ \\ \cline{1-4}
        $0$ & $241$ & $0$ & $38$ \\ \cline{1-4}
      \hline
    \end{tabular}
    }
  }
  \hfill
  \subfloat[Test data G]{
    \label{nnG}
    \scalebox{0.67}{
      \begin{tabular}{|c|c|c|c|}
      \hline
        $270$ & $7$ & $2$ & $0$ \\ \cline{1-4}
        $0$ & $280$ & $0$ & $0$ \\ \cline{1-4}
        $0$ & $39$ & $56$ & $0$ \\ \cline{1-4}
        $0$ & $255$ & $0$ & $24$ \\ \cline{1-4}
      \hline
    \end{tabular}
    }
  }
  \hfill
  \subfloat[Test data H]{
    \label{nnH}
    \scalebox{0.67}{
      \begin{tabular}{|c|c|c|c|}
      \hline
        $279$ & $0$ & $0$ & $0$ \\ \cline{1-4}
        $0$ & $274$ & $1$ & $2$ \\ \cline{1-4}
        $1$ & $15$ & $79$ & $0$ \\ \cline{1-4}
        $0$ & $254$ & $0$ & $25$ \\ \cline{1-4}
      \hline
    \end{tabular}
    }
  }
  \hfill
  \subfloat[Test data I]{
    \label{nnI}
    \scalebox{0.67}{
      \begin{tabular}{|c|c|c|c|}
      \hline
        $279$ & $0$ & $0$ & $0$ \\ \cline{1-4}
        $0$ & $274$ & $3$ & $2$ \\ \cline{1-4}
        $0$ & $42$ & $55$ & $0$ \\ \cline{1-4}
        $0$ & $179$ & $0$ & $100$ \\ \cline{1-4}
      \hline
    \end{tabular}
    }
  }
  \hfill
  \subfloat[Test data J]{
    \label{nnJ}
    \scalebox{0.67}{
      \begin{tabular}{|c|c|c|c|}
      \hline
        $264$ & $0$ & $0$ & $0$ \\ \cline{1-4}
        $0$ & $266$ & $0$ & $0$ \\ \cline{1-4}
        $0$ & $9$ & $73$ & $0$ \\ \cline{1-4}
        $81$ & $153$ & $0$ & $30$ \\ \cline{1-4}
      \hline
    \end{tabular}
    }
  }
  \hfill
  \subfloat[Test data K]{
    \label{nnK}
    \scalebox{0.67}{
      \begin{tabular}{|c|c|c|c|}
      \hline
        $268$ & $0$ & $0$ & $0$ \\ \cline{1-4}
        $0$ & $265$ & $0$ & $1$ \\ \cline{1-4}
        $0$ & $9$ & $75$ & $0$ \\ \cline{1-4}
        $0$ & $218$ & $0$ & $49$ \\ \cline{1-4}
      \hline
    \end{tabular}
    }
  }
  \hfill
  \subfloat[Test data L]{
    \label{nnL}
    \scalebox{0.67}{
      \begin{tabular}{|c|c|c|c|}
      \hline
        $267$ & $0$ & $0$ & $0$ \\ \cline{1-4}
        $0$ & $267$ & $0$ & $0$ \\ \cline{1-4}
        $0$ & $3$ & $81$ & $0$ \\ \cline{1-4}
        $0$ & $240$ & $0$ & $26$ \\ \cline{1-4}
      \hline
    \end{tabular}
    }
  }
  \hfill
  \subfloat[Test data M]{
    \label{nnM}
    \scalebox{0.67}{
      \begin{tabular}{|c|c|c|c|}
      \hline
        $257$ & $0$ & $1$ & $9$ \\ \cline{1-4}
        $0$ & $0$ & $0$ & $267$ \\ \cline{1-4}
        $0$ & $0$ & $0$ & $86$ \\ \cline{1-4}
        $0$ & $0$ & $0$ & $265$ \\ \cline{1-4}
      \hline
    \end{tabular}
    }
  }
  \hfill
  \subfloat[Test data N]{
    \label{nnN}
    \scalebox{0.67}{
      \begin{tabular}{|c|c|c|c|}
      \hline
        $265$ & $0$ & $0$ & $0$ \\ \cline{1-4}
        $0$ & $62$ & $0$ & $206$ \\ \cline{1-4}
        $0$ & $0$ & $0$ & $86$ \\ \cline{1-4}
        $0$ & $0$ & $0$ & $267$ \\ \cline{1-4}
      \hline
    \end{tabular}
    }
  }
  \hfill
  \subfloat[Test data O]{
    \label{nnO}
    \scalebox{0.67}{
      \begin{tabular}{|c|c|c|c|}
      \hline
        $267$ & $0$ & $0$ & $0$ \\ \cline{1-4}
        $0$ & $135$ & $0$ & $131$ \\ \cline{1-4}
        $0$ & $0$ & $85$ & $0$ \\ \cline{1-4}
        $0$ & $0$ & $0$ & $265$ \\ \cline{1-4}
      \hline
    \end{tabular}
    }
  }
  \hfill
  \subfloat[Test data P]{
    \label{nnP}
    \scalebox{0.67}{
      \begin{tabular}{|c|c|c|c|}
      \hline
        $267$ & $1$ & $0$ & $0$ \\ \cline{1-4}
        $1$ & $205$ & $0$ & $60$ \\ \cline{1-4}
        $0$ & $0$ & $83$ & $0$ \\ \cline{1-4}
        $2$ & $123$ & $0$ & $142$ \\ \cline{1-4}
      \hline
    \end{tabular}
    }
  }
  \hfill
  \subfloat[Test data Q]{
    \label{nnQ}
    \scalebox{0.67}{
      \begin{tabular}{|c|c|c|c|}
      \hline
        $266$ & $0$ & $0$ & $0$ \\ \cline{1-4}
        $0$ & $1$ & $0$ & $265$ \\ \cline{1-4}
        $0$ & $0$ & $86$ & $0$ \\ \cline{1-4}
        $1$ & $2$ & $0$ & $262$ \\ \cline{1-4}
      \hline
    \end{tabular}
    }
  }
  \hfill
  \subfloat[Test data R]{
    \label{nnR}
    \scalebox{0.67}{
      \begin{tabular}{|c|c|c|c|}
      \hline
        $197$ & $44$ & $26$ & $0$ \\ \cline{1-4}
        $0$ & $265$ & $1$ & $0$ \\ \cline{1-4}
        $0$ & $13$ & $72$ & $0$ \\ \cline{1-4}
        $0$ & $0$ & $0$ & $266$ \\ \cline{1-4}
      \hline
    \end{tabular}
    }
  }
  \hfill
  \subfloat[Test data S]{
    \label{nnS}
    \scalebox{0.67}{
      \begin{tabular}{|c|c|c|c|}
      \hline
        $246$ & $13$ & $6$ & $0$ \\ \cline{1-4}
        $0$ & $266$ & $0$ & $0$ \\ \cline{1-4}
        $0$ & $2$ & $83$ & $0$ \\ \cline{1-4}
        $8$ & $255$ & $0$ & $4$ \\ \cline{1-4}
      \hline
    \end{tabular}
    }
  }
  \hfill
  \subfloat[Test data T]{
    \label{nnT}
    \scalebox{0.67}{
      \begin{tabular}{|c|c|c|c|}
      \hline
        $251$ & $16$ & $0$ & $0$ \\ \cline{1-4}
        $0$ & $253$ & $0$ & $14$ \\ \cline{1-4}
        $0$ & $0$ & $85$ & $0$ \\ \cline{1-4}
        $9$ & $25$ & $0$ & $232$ \\ \cline{1-4}
      \hline
    \end{tabular}
    }
  }
  \hfill
  \subfloat[Test data U]{
    \label{nnU}
    \scalebox{0.67}{
      \begin{tabular}{|c|c|c|c|}
      \hline
        $260$ & $5$ & $2$ & $0$ \\ \cline{1-4}
        $0$ & $263$ & $3$ & $0$ \\ \cline{1-4}
        $1$ & $1$ & $82$ & $0$ \\ \cline{1-4}
        $163$ & $95$ & $0$ & $9$ \\ \cline{1-4}
      \hline
    \end{tabular}
    }
  }
  \hfill
  \subfloat[Test data V]{
    \label{nnV}
    \scalebox{0.67}{
      \begin{tabular}{|c|c|c|c|}
      \hline
        $266$ & $0$ & $0$ & $0$ \\ \cline{1-4}
        $1$ & $241$ & $26$ & $0$ \\ \cline{1-4}
        $0$ & $47$ & $37$ & $0$ \\ \cline{1-4}
        $2$ & $262$ & $1$ & $1$ \\ \cline{1-4}
      \hline
    \end{tabular}
    }
  }
  \hfill
  \subfloat[Test data W]{
    \label{nnW}
    \scalebox{0.67}{
      \begin{tabular}{|c|c|c|c|}
      \hline
        $231$ & $11$ & $24$ & $0$ \\ \cline{1-4}
        $0$ & $245$ & $21$ & $0$ \\ \cline{1-4}
        $0$ & $11$ & $75$ & $0$ \\ \cline{1-4}
        $0$ & $266$ & $0$ & $0$ \\ \cline{1-4}
      \hline
    \end{tabular}
    }
  }
  \hfill
  \subfloat[Test data X]{
    \label{nnX}
    \scalebox{0.67}{
      \begin{tabular}{|c|c|c|c|}
      \hline
        $233$ & $13$ & $21$ & $0$ \\ \cline{1-4}
        $0$ & $225$ & $42$ & $0$ \\ \cline{1-4}
        $0$ & $1$ & $83$ & $0$ \\ \cline{1-4}
        $0$ & $258$ & $9$ & $0$ \\ \cline{1-4}
      \hline
    \end{tabular}
    }
  }
  \hfill
  \subfloat[Test data Y]{
    \label{nnY}
    \scalebox{0.67}{
      \begin{tabular}{|c|c|c|c|}
      \hline
        $267$ & $0$ & $0$ & $0$ \\ \cline{1-4}
        $0$ & $202$ & $65$ & $0$ \\ \cline{1-4}
        $0$ & $19$ & $66$ & $0$ \\ \cline{1-4}
        $8$ & $258$ & $1$ & $0$ \\ \cline{1-4}
      \hline
    \end{tabular}
    }
  }
  \hfill
  \subfloat[Test data Z]{
    \label{nnZ}
    \scalebox{0.67}{
      \begin{tabular}{|c|c|c|c|}
      \hline
        $266$ & $0$ & $0$ & $0$ \\ \cline{1-4}
        $58$ & $198$ & $10$ & $0$ \\ \cline{1-4}
        $0$ & $4$ & $80$ & $0$ \\ \cline{1-4}
        $1$ & $266$ & $0$ & $0$ \\ \cline{1-4}
      \hline
    \end{tabular}
    }
  }
  \hfill
\setcounter{subtable}{0}
  \hfill
  \subfloat[Test data AA]{
    \label{nnAA}
    \scalebox{0.67}{
      \begin{tabular}{|c|c|c|c|}
      \hline
        $253$ & $9$ & $4$ & $0$ \\ \cline{1-4}
        $0$ & $218$ & $49$ & $0$ \\ \cline{1-4}
        $0$ & $4$ & $81$ & $0$ \\ \cline{1-4}
        $3$ & $240$ & $8$ & $14$ \\ \cline{1-4}
      \hline
    \end{tabular}
    }
  }
  \hfill
  \subfloat[Test data AB]{
    \label{nnAB}
    \scalebox{0.67}{
      \begin{tabular}{|c|c|c|c|}
      \hline
        $262$ & $2$ & $2$ & $0$ \\ \cline{1-4}
        $0$ & $263$ & $3$ & $0$ \\ \cline{1-4}
        $0$ & $0$ & $84$ & $0$ \\ \cline{1-4}
        $56$ & $195$ & $3$ & $12$ \\ \cline{1-4}
      \hline
    \end{tabular}
    }
  }
\end{table}

\begin{table}[b]
  \centering
  \caption{Matrix of test results with normals and k-means. Diagonal toward lower right corner is correct. Vertical columns represent number of labels inferred as mouse operation, pad operation, sitting still, and typing, respectively, from left to right. Horizontal rows represent, from top, number of labels where actual label is mouse operation, pad operation, sitting still, or typing.}
  \label{normalkmeansmatrix}
  \subfloat[Test data A]{
    \label{nkA}
    \scalebox{0.67}{
      \begin{tabular}{|c|c|c|c|}
      \hline
        $268$ & $5$ & $0$ & $0$ \\ \cline{1-4}
        $73$ & $197$ & $1$ & $3$ \\ \cline{1-4}
        $7$ & $5$ & $84$ & $0$ \\ \cline{1-4}
        $16$ & $253$ & $3$ & $4$ \\ \cline{1-4}
      \hline
    \end{tabular}
    }
  }
  \hfill
  \subfloat[Test data B]{
    \label{nkB}
    \scalebox{0.67}{
      \begin{tabular}{|c|c|c|c|}
      \hline
        $228$ & $30$ & $12$ & $7$ \\ \cline{1-4}
        $1$ & $254$ & $3$ & $17$ \\ \cline{1-4}
        $12$ & $3$ & $78$ & $0$ \\ \cline{1-4}
        $3$ & $17$ & $1$ & $255$ \\ \cline{1-4}
      \hline
    \end{tabular}
    }
  }
  \hfill
  \subfloat[Test data C]{
    \label{nkC}
    \scalebox{0.67}{
      \begin{tabular}{|c|c|c|c|}
      \hline
        $279$ & $0$ & $0$ & $0$ \\ \cline{1-4}
        $0$ & $0$ & $0$ & $278$ \\ \cline{1-4}
        $0$ & $3$ & $91$ & $0$ \\ \cline{1-4}
        $0$ & $0$ & $0$ & $276$ \\ \cline{1-4}
      \hline
    \end{tabular}
    }
  }
  \hfill
  \subfloat[Test data D]{
    \label{nkD}
    \scalebox{0.67}{
      \begin{tabular}{|c|c|c|c|}
      \hline
        $216$ & $9$ & $17$ & $35$ \\ \cline{1-4}
        $8$ & $57$ & $30$ & $184$ \\ \cline{1-4}
        $9$ & $0$ & $86$ & $0$ \\ \cline{1-4}
        $17$ & $204$ & $29$ & $25$ \\ \cline{1-4}
      \hline
    \end{tabular}
    }
  }
  \hfill
  \subfloat[Test data E]{
    \label{nkE}
    \scalebox{0.67}{
      \begin{tabular}{|c|c|c|c|}
      \hline
        $88$ & $0$ & $0$ & $190$ \\ \cline{1-4}
        $0$ & $273$ & $4$ & $0$ \\ \cline{1-4}
        $0$ & $38$ & $59$ & $0$ \\ \cline{1-4}
        $2$ & $267$ & $4$ & $4$ \\ \cline{1-4}
      \hline
    \end{tabular}
    }
  }
  \hfill
  \subfloat[Test data F]{
    \label{nkF}
    \scalebox{0.67}{
      \begin{tabular}{|c|c|c|c|}
      \hline
        $117$ & $16$ & $147$ & $0$ \\ \cline{1-4}
        $0$ & $263$ & $14$ & $2$ \\ \cline{1-4}
        $0$ & $0$ & $97$ & $0$ \\ \cline{1-4}
        $1$ & $248$ & $4$ & $26$ \\ \cline{1-4}
      \hline
    \end{tabular}
    }
  }
  \hfill
  \subfloat[Test data G]{
    \label{nkG}
    \scalebox{0.67}{
      \begin{tabular}{|c|c|c|c|}
      \hline
        $39$ & $9$ & $231$ & $0$ \\ \cline{1-4}
        $11$ & $232$ & $32$ & $5$ \\ \cline{1-4}
        $0$ & $11$ & $84$ & $0$ \\ \cline{1-4}
        $2$ & $256$ & $2$ & $19$ \\ \cline{1-4}
      \hline
    \end{tabular}
    }
  }
  \hfill
  \subfloat[Test data H]{
    \label{nkH}
    \scalebox{0.67}{
      \begin{tabular}{|c|c|c|c|}
      \hline
        $261$ & $0$ & $0$ & $18$ \\ \cline{1-4}
        $0$ & $271$ & $3$ & $3$ \\ \cline{1-4}
        $0$ & $2$ & $93$ & $0$ \\ \cline{1-4}
        $1$ & $264$ & $6$ & $8$ \\ \cline{1-4}
      \hline
    \end{tabular}
    }
  }
  \hfill
  \subfloat[Test data I]{
    \label{nkI}
    \scalebox{0.67}{
      \begin{tabular}{|c|c|c|c|}
      \hline
        $182$ & $1$ & $95$ & $1$ \\ \cline{1-4}
        $0$ & $275$ & $4$ & $0$ \\ \cline{1-4}
        $1$ & $6$ & $90$ & $0$ \\ \cline{1-4}
        $2$ & $231$ & $9$ & $37$ \\ \cline{1-4}
      \hline
    \end{tabular}
    }
  }
  \hfill
  \subfloat[Test data J]{
    \label{nkJ}
    \scalebox{0.67}{
      \begin{tabular}{|c|c|c|c|}
      \hline
        $264$ & $0$ & $0$ & $0$ \\ \cline{1-4}
        $3$ & $248$ & $0$ & $15$ \\ \cline{1-4}
        $1$ & $62$ & $19$ & $0$ \\ \cline{1-4}
        $8$ & $14$ & $0$ & $242$ \\ \cline{1-4}
      \hline
    \end{tabular}
    }
  }
  \hfill
  \subfloat[Test data K]{
    \label{nkK}
    \scalebox{0.67}{
      \begin{tabular}{|c|c|c|c|}
      \hline
        $253$ & $14$ & $0$ & $1$ \\ \cline{1-4}
        $0$ & $258$ & $2$ & $6$ \\ \cline{1-4}
        $0$ & $84$ & $0$ & $0$ \\ \cline{1-4}
        $0$ & $6$ & $0$ & $261$ \\ \cline{1-4}
      \hline
    \end{tabular}
    }
  }
  \hfill
  \subfloat[Test data L]{
    \label{nkL}
    \scalebox{0.67}{
      \begin{tabular}{|c|c|c|c|}
      \hline
        $267$ & $0$ & $0$ & $0$ \\ \cline{1-4}
        $0$ & $267$ & $0$ & $0$ \\ \cline{1-4}
        $0$ & $57$ & $27$ & $0$ \\ \cline{1-4}
        $2$ & $26$ & $0$ & $238$ \\ \cline{1-4}
      \hline
    \end{tabular}
    }
  }
  \hfill
  \subfloat[Test data M]{
    \label{nkM}
    \scalebox{0.67}{
      \begin{tabular}{|c|c|c|c|}
      \hline
        $159$ & $26$ & $1$ & $81$ \\ \cline{1-4}
        $46$ & $21$ & $22$ & $178$ \\ \cline{1-4}
        $1$ & $0$ & $5$ & $80$ \\ \cline{1-4}
        $33$ & $7$ & $9$ & $216$ \\ \cline{1-4}
      \hline
    \end{tabular}
    }
  }
  \hfill
  \subfloat[Test data N]{
    \label{nkN}
    \scalebox{0.67}{
      \begin{tabular}{|c|c|c|c|}
      \hline
        $265$ & $0$ & $0$ & $0$ \\ \cline{1-4}
        $0$ & $14$ & $1$ & $253$ \\ \cline{1-4}
        $0$ & $1$ & $0$ & $85$ \\ \cline{1-4}
        $0$ & $0$ & $0$ & $267$ \\ \cline{1-4}
      \hline
    \end{tabular}
    }
  }
  \hfill
  \subfloat[Test data O]{
    \label{nkO}
    \scalebox{0.67}{
      \begin{tabular}{|c|c|c|c|}
      \hline
        $252$ & $5$ & $1$ & $9$ \\ \cline{1-4}
        $7$ & $31$ & $4$ & $224$ \\ \cline{1-4}
        $0$ & $1$ & $84$ & $0$ \\ \cline{1-4}
        $1$ & $3$ & $0$ & $261$ \\ \cline{1-4}
      \hline
    \end{tabular}
    }
  }
  \hfill
  \subfloat[Test data P]{
    \label{nkP}
    \scalebox{0.67}{
      \begin{tabular}{|c|c|c|c|}
      \hline
        $268$ & $0$ & $0$ & $0$ \\ \cline{1-4}
        $49$ & $46$ & $3$ & $168$ \\ \cline{1-4}
        $0$ & $22$ & $61$ & $0$ \\ \cline{1-4}
        $4$ & $16$ & $2$ & $245$ \\ \cline{1-4}
      \hline
    \end{tabular}
    }
  }
  \hfill
  \subfloat[Test data Q]{
    \label{nkQ}
    \scalebox{0.67}{
      \begin{tabular}{|c|c|c|c|}
      \hline
        $248$ & $15$ & $0$ & $3$ \\ \cline{1-4}
        $0$ & $13$ & $3$ & $250$ \\ \cline{1-4}
        $0$ & $2$ & $84$ & $0$ \\ \cline{1-4}
        $2$ & $8$ & $2$ & $253$ \\ \cline{1-4}
      \hline
    \end{tabular}
    }
  }
  \hfill
  \subfloat[Test data R]{
    \label{nkR}
    \scalebox{0.67}{
      \begin{tabular}{|c|c|c|c|}
      \hline
        $237$ & $3$ & $27$ & $0$ \\ \cline{1-4}
        $1$ & $254$ & $0$ & $11$ \\ \cline{1-4}
        $0$ & $10$ & $75$ & $0$ \\ \cline{1-4}
        $2$ & $14$ & $5$ & $245$ \\ \cline{1-4}
      \hline
    \end{tabular}
    }
  }
  \hfill
  \subfloat[Test data S]{
    \label{nkS}
    \scalebox{0.67}{
      \begin{tabular}{|c|c|c|c|}
      \hline
        $210$ & $33$ & $22$ & $0$ \\ \cline{1-4}
        $0$ & $266$ & $0$ & $0$ \\ \cline{1-4}
        $0$ & $8$ & $77$ & $0$ \\ \cline{1-4}
        $10$ & $124$ & $1$ & $132$ \\ \cline{1-4}
      \hline
    \end{tabular}
    }
  }
  \hfill
  \subfloat[Test data T]{
    \label{nkT}
    \scalebox{0.67}{
      \begin{tabular}{|c|c|c|c|}
      \hline
        $262$ & $4$ & $1$ & $0$ \\ \cline{1-4}
        $0$ & $30$ & $0$ & $237$ \\ \cline{1-4}
        $0$ & $1$ & $84$ & $0$ \\ \cline{1-4}
        $8$ & $11$ & $1$ & $246$ \\ \cline{1-4}
      \hline
    \end{tabular}
    }
  }
  \hfill
  \subfloat[Test data U]{
    \label{nkU}
    \scalebox{0.67}{
      \begin{tabular}{|c|c|c|c|}
      \hline
        $264$ & $0$ & $3$ & $0$ \\ \cline{1-4}
        $2$ & $255$ & $0$ & $9$ \\ \cline{1-4}
        $0$ & $3$ & $81$ & $0$ \\ \cline{1-4}
        $108$ & $8$ & $0$ & $151$ \\ \cline{1-4}
      \hline
    \end{tabular}
    }
  }
  \hfill
  \subfloat[Test data V]{
    \label{nkV}
    \scalebox{0.67}{
      \begin{tabular}{|c|c|c|c|}
      \hline
        $266$ & $0$ & $0$ & $0$ \\ \cline{1-4}
        $0$ & $266$ & $2$ & $0$ \\ \cline{1-4}
        $0$ & $48$ & $36$ & $0$ \\ \cline{1-4}
        $0$ & $113$ & $0$ & $153$ \\ \cline{1-4}
      \hline
    \end{tabular}
    }
  }
  \hfill
  \subfloat[Test data W]{
    \label{nkW}
    \scalebox{0.67}{
      \begin{tabular}{|c|c|c|c|}
      \hline
        $230$ & $5$ & $31$ & $0$ \\ \cline{1-4}
        $1$ & $265$ & $0$ & $0$ \\ \cline{1-4}
        $0$ & $7$ & $79$ & $0$ \\ \cline{1-4}
        $1$ & $262$ & $0$ & $3$ \\ \cline{1-4}
      \hline
    \end{tabular}
    }
  }
  \hfill
  \subfloat[Test data X]{
    \label{nkX}
    \scalebox{0.67}{
      \begin{tabular}{|c|c|c|c|}
      \hline
        $84$ & $97$ & $86$ & $0$ \\ \cline{1-4}
        $0$ & $260$ & $7$ & $0$ \\ \cline{1-4}
        $0$ & $19$ & $65$ & $0$ \\ \cline{1-4}
        $8$ & $247$ & $2$ & $10$ \\ \cline{1-4}
      \hline
    \end{tabular}
    }
  }
  \hfill
  \subfloat[Test data Y]{
    \label{nkY}
    \scalebox{0.67}{
      \begin{tabular}{|c|c|c|c|}
      \hline
        $217$ & $31$ & $19$ & $0$ \\ \cline{1-4}
        $0$ & $264$ & $3$ & $0$ \\ \cline{1-4}
        $0$ & $31$ & $54$ & $0$ \\ \cline{1-4}
        $3$ & $230$ & $1$ & $33$ \\ \cline{1-4}
      \hline
    \end{tabular}
    }
  }
  \hfill
  \subfloat[Test data Z]{
    \label{nkZ}
    \scalebox{0.67}{
      \begin{tabular}{|c|c|c|c|}
      \hline
        $251$ & $8$ & $5$ & $2$ \\ \cline{1-4}
        $51$ & $210$ & $0$ & $5$ \\ \cline{1-4}
        $0$ & $22$ & $62$ & $0$ \\ \cline{1-4}
        $0$ & $189$ & $1$ & $77$ \\ \cline{1-4}
      \hline
    \end{tabular}
    }
  }
  \hfill
\setcounter{subtable}{0}
  \hfill
  \subfloat[Test data AA]{
    \label{nkAA}
    \scalebox{0.67}{
      \begin{tabular}{|c|c|c|c|}
      \hline
        $263$ & $3$ & $0$ & $0$ \\ \cline{1-4}
        $2$ & $263$ & $2$ & $0$ \\ \cline{1-4}
        $0$ & $13$ & $72$ & $0$ \\ \cline{1-4}
        $0$ & $29$ & $0$ & $236$ \\ \cline{1-4}
      \hline
    \end{tabular}
    }
  }
  \hfill
  \subfloat[Test data AB]{
    \label{nkAB}
    \scalebox{0.67}{
      \begin{tabular}{|c|c|c|c|}
      \hline
        $266$ & $0$ & $0$ & $0$ \\ \cline{1-4}
        $4$ & $260$ & $0$ & $2$ \\ \cline{1-4}
        $0$ & $2$ & $82$ & $0$ \\ \cline{1-4}
        $24$ & $2$ & $1$ & $239$ \\ \cline{1-4}
      \hline
    \end{tabular}
    }
  }
\end{table}


}